\documentclass{beamer}
\usepackage[utf8]{inputenc}

\usepackage{utopia} %font utopia imported

\usetheme{Madrid}
\usecolortheme{default}

%------------------------------------------------------------
%This block of code defines the information to appear in the
%Title page
\title[Operating Systems] %optional
{Operating Systems}

\subtitle{Beamer version 0.1}

\author[Fabricio Bortoluzzi] % (optional)
{Fabricio Bortoluzzi\inst{1}}

\institute[UNIVALI] % (optional)
{
  \inst{1}%
  Computer Networks Laboratory\\
  in collaboration with\\
  Laboratory of Embedded and Distributed Systems\\
}

\date[UNIVALI 2019] % (optional)
{www.univali.br\\leds.acad.univali.br}

%\logo{\includegraphics[height=1.5cm]{univali-logo.jpg}}

%End of title page configuration block
%------------------------------------------------------------



%------------------------------------------------------------
%The next block of commands puts the table of contents at the 
%beginning of each section and highlights the current section:

\AtBeginSection[]
{
  \begin{frame}
    \frametitle{Table of Contents}
    \tableofcontents[currentsection]
  \end{frame}
}
%------------------------------------------------------------


\begin{document}

%The next statement creates the title page.
\frame{\titlepage}


%---------------------------------------------------------
%This block of code is for the table of contents after
%the title page
\begin{frame}
\frametitle{Table of Contents}
\tableofcontents
\end{frame}
%---------------------------------------------------------


\section{Introduction}

%---------------------------------------------------------
%Basic frame
\begin{frame}
    \frametitle{Hardware review}
    
    \begin{itemize}
        \item 
        A modern  computer  consists  of  one  or  more
        processors, some main  memory, disks,  printers,  
        a  keyboard,  a  mouse, a  display,  
        network  interfaces,  and  various other 
        input/output devices.

        \item 
        All in all, a complex system.
        
        \item 
        If every application programmer had to understand
        how all these things work in detail,
        no code would ever get written.

    \end{itemize}    

\end{frame}
%---------------------------------------------------------

%---------------------------------------------------------
%Basic frame
\begin{frame}
    \frametitle{Wider definition of an Operating Systems}
    
    The operating system (OS) is a layer inserted 
    between the hardware/software interface in order
    to provide a 
    \begin{itemize}
        \item better
        \item simpler
        \item cleaner
    \end{itemize}, 
    model of the computer.

\end{frame}
%---------------------------------------------------------

%---------------------------------------------------------
%Basic frame
\begin{frame}
    \frametitle{The shell is not the OS}
    
    Users interact with the shell, either via 
    \begin{itemize}
        \item a Text Console; and/or
        \item the GUI - Graphical User Interface.
    \end{itemize}

\end{frame}
%---------------------------------------------------------

\section{Memory Management}

%---------------------------------------------------------
%Highlighting text
\begin{frame}
    \frametitle{First Frame of Memory Management}
    
    In this slide, some important text will be
    \alert{highlighted} because it's important.
    Please, don't abuse it.
    
    \begin{block}{Remark}
    Sample text
    \end{block}
    
    \begin{alertblock}{Important theorem}
    Sample text
    \end{alertblock}
    
    \begin{examples}
    Sample text
    \end{examples}
    \end{frame}
    %---------------------------------------------------------
    
\section{File Systems}

%---------------------------------------------------------
%Basic frame
\begin{frame}
    \frametitle{First Frame of File Systems}
    
    Text of the first frame of file systems.

\end{frame}
%---------------------------------------------------------

\section{Input/Output}

%---------------------------------------------------------
%Basic frame
\begin{frame}
    \frametitle{First Frame of I/O}
    
    Text of the first frame of I/O.

\end{frame}
%---------------------------------------------------------

\section{Deadlocks}

%---------------------------------------------------------
%Basic frame
\begin{frame}
    \frametitle{Deadlocks}
    
    Text of the first frame of Deadlocks.

\end{frame}
%---------------------------------------------------------


%\section{Examples}

%---------------------------------------------------------
%Changing visivility of the text
\begin{frame}
    \frametitle{Sample frame title}
    This is a text in second frame. For the sake of showing an example.
    
    \begin{itemize}
        \item<1-> Text visible on slide 1
        \item<2-> Text visible on slide 2
        \item<3> Text visible on slides 3
        \item<4-> Text visible on slide 4
    \end{itemize}
    \end{frame}
    
    %---------------------------------------------------------
    
    
    %---------------------------------------------------------
    %Example of the \pause command
    \begin{frame}
    In this slide \pause
    
    the text will be partially visible \pause
    
    And finally everything will be there
    \end{frame}
    %---------------------------------------------------------
    
    \section{Processes and Threads}
    
    %---------------------------------------------------------
    %Highlighting text
    \begin{frame}
    \frametitle{Sample frame title}
    
    In this slide, some important text will be
    \alert{highlighted} because it's important.
    Please, don't abuse it.
    
    \begin{block}{Remark}
    Sample text
    \end{block}
    
    \begin{alertblock}{Important theorem}
    Sample text in red box
    \end{alertblock}
    
    \begin{examples}
    Sample text in green box. The title of the block is ``Examples".
    \end{examples}
    \end{frame}
    %---------------------------------------------------------
    
    
    %---------------------------------------------------------
    %Two columns
    \begin{frame}
    \frametitle{Two-column slide}
    
    \begin{columns}
    
    \column{0.5\textwidth}
    This is a text in first column.
    $$E=mc^2$$
    \begin{itemize}
    \item First item
    \item Second item
    \end{itemize}
    
    \column{0.5\textwidth}
    This text will be in the second column
    and on a second tought this is a nice looking
    layout in some cases.
    \end{columns}
    \end{frame}
    %---------------------------------------------------------
    
    

\end{document}